% =============================================================
% TABLAS COMPARATIVAS — ESTADO DEL ARTE  (versión 2 — verificada)
% Correcciones respecto a v1:
%   (1) MM-Skin: Acc 95.63 % → SkinVL-PubMM (no SkinVL-MM).
%       SkinVL-MM tiene Acc 95.27 % en Patch16. Se incluyen ambas.
%   (2) AutomatedReport: se agrega ROUGE-L = 0.3828 (antes "—").
%   (3) MPXV-CNN: se indica explícitamente que AUC proviene del
%       cohorte externo y Sens./Espec. del cohorte de validación.
%   (4) Ensemble DL: nota aclaratoria sobre métricas de clase NV.
%
% REQUISITOS EN EL PREÁMBULO:
%   \usepackage{booktabs}
%   \usepackage{graphicx}
%   \usepackage{xcolor}
%   \usepackage{colortbl}
% =============================================================

\definecolor{thiswork}{RGB}{230, 242, 255}

% ──────────────────────────────────────────────────────────────
%  TABLA 1 — Clasificación de lesiones dermatológicas
% ──────────────────────────────────────────────────────────────
\begin{table}[htbp]
\centering
\caption{%
  Comparativa del estado del arte en clasificación de lesiones dermatológicas.
  $^\dagger$~indica arquitectura basada en VLM\@.
  \textit{FT}: ajuste fino supervisado;
  \textit{ZS}: evaluación \textit{zero-shot}.
  \textit{Cls.}: número de clases objetivo.
  ``—'' indica métrica no reportada.
  Las filas sombreadas corresponden a los modelos propuestos en este trabajo
  y serán completadas con los resultados experimentales.
}
\label{tab:sota_clasificacion}
\resizebox{\textwidth}{!}{%
\begin{tabular}{@{}lllccrrrrr@{}}
\toprule
\textbf{Referencia}
  & \textbf{Arquitectura}
  & \textbf{Dataset}
  & \textbf{Modo}
  & \textbf{Cls.}
  & \textbf{Acc. (\%)}
  & \textbf{AUC-ROC}
  & \textbf{F1}
  & \textbf{Sens.}
  & \textbf{Espec.} \\
\midrule

% ── CNN / ViT ────────────────────────────────────────────────
\multicolumn{10}{@{}l}{\small\itshape Enfoques basados en CNN / ViT} \\[2pt]

\cite{SkinHealthMate2024}
  & EfficientNetB1$+$B5 (ens.)
  & HAM10000
  & FT & 7
  & 93.00 & — & $^{(\alpha)}$ & — & — \\

\cite{SkinSavvy22025}
  & EfficientNetV2~$+$~Swin-Transformer
  & HAM10000
  & FT & 7
  & 77.70 & — & — & — & — \\

\cite{TwoPath2024}
  & EfficientNet~$+$~MobileNetV2
  & ISIC-2019
  & FT & 8
  & 70.00 & — & 0.667 & 0.750 & — \\

\cite{EnsembleDL2025}
  & EfficientNetB3$+$ResNet50$+$DenseNet121
  & ISIC (bin.)
  & FT & 2
  & 97.67 & — & $^{(\beta)}$ & $^{(\beta)}$ & — \\

\cite{MPXVCNN2023}
  & ResNet34
  & Multi-dataset (mpox)
  & FT & 2
  & — & $0.966^{(\gamma)}$ & — & $0.830^{(\gamma)}$ & $0.965^{(\gamma)}$ \\

\cite{PanDerm2024}
  & PanDerm (ViT-Large)
  & 2.1M imág.~/ 11 inst.
  & FT & bin.
  & — & \textbf{0.964} & — & 0.893 & — \\

\addlinespace[5pt]
% ── VLM ──────────────────────────────────────────────────────
\multicolumn{10}{@{}l}{\small\itshape Enfoques basados en VLM$\,^\dagger$} \\[2pt]

\cite{MpoxVLM2024}$^\dagger$
  & CLIP~$+$~ViT~$+$~LLaMA-2-7B
  & Mpox dataset
  & FT & 2
  & 90.38 & 0.9516 & 0.857 & 0.844 & — \\

\cite{DermLIP2025}$^\dagger$
  & DermLIP (CLIP ViT-B16/GPT77)
  & Derm1M
  & ZS & múlt.
  & 56.07 & 0.731$^{(\delta)}$ & — & — & — \\

\cite{MMSkin2025}$^\dagger$
  & SkinVL-MM (LLaVA-Med-v1.5-7B~$+$~LoRA)
  & MM-Skin
  & FT & $^{(\varepsilon)}$
  & 95.27 & — & — & — & — \\

\cite{MMSkin2025}$^\dagger$
  & SkinVL-PubMM (LLaVA-Med-v1.5-7B~$+$~LoRA)
  & MM-Skin~$+$~públicos
  & FT & $^{(\varepsilon)}$
  & \textbf{95.63} & — & — & — & — \\

\cite{DermatoLlama2025}$^\dagger$
  & DermatoLlama (Eval2~MCQA)
  & DermaSynth (sint.)
  & FT & múlt.
  & 91.00 & — & — & — & — \\

\midrule
% ── Este trabajo ─────────────────────────────────────────────
\multicolumn{10}{@{}l}{\small\itshape Este trabajo} \\[2pt]

\rowcolor{thiswork}
(a) & DenseNet121
  & HAM10000$^{(\zeta)}$
  & FT & 7
  & — & — & — & — & — \\

\rowcolor{thiswork}
(b) & ViT-B/16
  & HAM10000$^{(\zeta)}$
  & FT & 7
  & — & — & — & — & — \\

\rowcolor{thiswork}
(a$'$) & ConvNeXt-Tiny
  & HAM10000$^{(\zeta)}$
  & FT & 7
  & — & — & — & — & — \\

\rowcolor{thiswork}
(b$'$) & Swin-Tiny
  & HAM10000$^{(\zeta)}$
  & FT & 7
  & — & — & — & — & — \\

\rowcolor{thiswork}
(c)$^\dagger$ & Mini-VLM (CLIP ViT-B/16~$+$~Phi-3-mini, 3.9B)
  & HAM10000~$+$~HIBA
  & FT & 7
  & — & — & — & — & — \\

\rowcolor{thiswork}
(d)$^\dagger$ & LLaVA-Med 7B~$+$~LoRA~$(r{=}64)$
  & HAM10000~$+$~HIBA
  & FT & 7
  & — & — & — & — & — \\

\bottomrule
\end{tabular}%
}
\vspace{4pt}
\begin{minipage}{\textwidth}
  \footnotesize
  $^{(\alpha)}$~F1 reportado por clase: NV\,=\,0.98, MEL\,=\,0.55, BCC\,=\,0.84;
  F1-macro\,=\,0.81, F1-weighted\,=\,0.93~\cite{SkinHealthMate2024}.\\
  $^{(\beta)}$~F1 = 0.976 y Sensibilidad = 0.963 corresponden a la clase NV
  (benigna); métricas globales del ensemble no reportadas en el original~\cite{EnsembleDL2025}.\\
  $^{(\gamma)}$~AUC\,=\,0.966 y Sensibilidad\,=\,0.91 provienen del cohorte de
  testeo externo; Sensibilidad\,=\,0.83 y Especificidad\,=\,0.965 del cohorte
  de validación cruzada. Los valores de Sens./Espec.\ mostrados en esta fila
  corresponden al cohorte de validación~\cite{MPXVCNN2023}.\\
  $^{(\delta)}$~AUROC reportado sobre el benchmark SkinCon (32 conceptos);
  Recall@10\,=\,59.98\,\% (imagen$\to$texto) en holdout set de Derm1M~\cite{DermLIP2025}.\\
  $^{(\varepsilon)}$~Exactitud a nivel de parche (\textit{patch-level}) en imágenes
  de patología (Patch16). SkinVL-MM entrenado únicamente con MM-Skin;
  SkinVL-PubMM entrenado con MM-Skin~$+$~9 datasets públicos~\cite{MMSkin2025}.
  No directamente comparable con clasificación de lesión completa sobre HAM10000.\\
  $^{(\zeta)}$~Filas (a)(b)(a$'$)(b$'$): entrenamiento sobre HAM10000; HIBA se
  utiliza únicamente como conjunto adicional de evaluación OOD (1\,616 imágenes
  mapeadas a las 7 clases HAM10000; el mapeo SCC$\to$AKIEC es aproximado
  porque AKIEC original sólo incluye carcinoma intraepitelial, no SCC invasivo).
  ConvNeXt-Tiny y Swin-Tiny incorporan el \textit{recipe} \texttt{modern}
  (RandAugment, RandomErasing, MixUp/CutMix, Focal Loss, label smoothing,
  warmup, two-stage fine-tuning con LLRD\,=\,0.75 y TTA por horizontal flip).\\
  La comparabilidad directa entre filas es limitada debido a diferencias
  en datasets, número de clases y protocolos de evaluación.
\end{minipage}
\end{table}


% ──────────────────────────────────────────────────────────────
%  TABLA 2 — Generación automática de reportes clínicos
% ──────────────────────────────────────────────────────────────
\begin{table}[htbp]
\centering
\caption{%
  Comparativa del estado del arte en generación automática de reportes
  clínicos dermatológicos mediante VLMs.
  BLEU-4 y ROUGE en escala 0--100 salvo indicación contraria;
  BERTScore F1 en escala 0--1.
  ``—'' indica métrica no reportada.
  Las filas sombreadas serán completadas con los resultados experimentales.
}
\label{tab:sota_reportes}
\resizebox{\textwidth}{!}{%
\begin{tabular}{@{}lllrrrrrl@{}}
\toprule
\textbf{Referencia}
  & \textbf{Arquitectura}
  & \textbf{Dataset (eval.)}
  & \textbf{BLEU-4}
  & \textbf{ROUGE-1}
  & \textbf{ROUGE-L}
  & \textbf{METEOR}
  & \textbf{BERTScore F1}
  & \textbf{Eval.\ clínica} \\
\midrule

\cite{AutomatedReport2025}
  & LLaVA-v1.5-13B (KD$^{(\zeta)}$)
  & Dermoscopia priv.\ (305 img.)
  & $55.59^{(\eta)}$
  & 54.38
  & 38.28
  & —
  & 0.9025
  & — \\

\cite{DermatoLlama2025}
  & DermatoLlama (Full~$+$~IF)
  & DERM12345 (Eval1)
  & $65.59^{(\theta)}$
  & 93.49$^{(\theta)}$
  & 93.44$^{(\theta)}$
  & —
  & —
  & Acc.\ MCQA: 90.93\% \\

\cite{MMSkin2025}
  & SkinVL-MM (LLaVA-Med-v1.5-7B~$+$~LoRA)
  & MM-Skin (patología)
  & 22.04$^{(\iota)}$
  & —
  & —
  & —
  & —
  & — \\

\cite{SkinGPT4_2024}
  & ViT~$+$~Q-Former~$+$~LLaMA-2-13B
  & Privado (52k~img.)
  & —
  & —
  & —
  & —
  & —
  & 80.63\%~correctos \\

\cite{PeFoMed2024}
  & ViT-EVA~$+$~LLaMA-2-7B~$+$~LoRA
  & IU-Xray$^{(\kappa)}$
  & —
  & —
  & 28.60
  & 15.70
  & —
  & WASM: 2.79/4 \\

\midrule
\multicolumn{9}{@{}l}{\small\itshape Este trabajo} \\[2pt]

\rowcolor{thiswork}
(c)$^\dagger$
  & Mini-VLM (CLIP ViT-B/16~$+$~Phi-3-mini, 3.9B)
  & MM-Skin (test split)
  & — & — & — & — & — & — \\

\rowcolor{thiswork}
(d)$^\dagger$
  & LLaVA-Med 7B~$+$~LoRA~$(r{=}64)$
  & MM-Skin (test split)
  & — & — & — & — & — & — \\

\bottomrule
\end{tabular}%
}
\vspace{4pt}
\begin{minipage}{\textwidth}
  \footnotesize
  $^{(\zeta)}$~Modelo estudiante entrenado mediante destilación de conocimiento
  (\textit{knowledge distillation}) con GPT-4V como modelo maestro~\cite{AutomatedReport2025}.\\
  $^{(\eta)}$~Reportado como SacreBLEU (escala 0--100); no equivalente a BLEU-4
  calculado con corpus BLEU estándar. ROUGE-1 y ROUGE-L en escala 0--100
  (originalmente reportados como 0.5438 y 0.3828 respectivamente).\\
  $^{(\theta)}$~Valores reportados en escala 0--1 en el original
  (BLEU-4\,=\,0.6559, ROUGE-1\,=\,0.9349, ROUGE-L\,=\,0.9344);
  aquí multiplicados por 100. Accuracy (Eval2 - PubMed MCQA)\,=\,0.9093.
  Los valores de ROUGE inusualmente altos pueden reflejar consistencia
  estructural en datos sintéticos~\cite{DermatoLlama2025}.\\
  $^{(\iota)}$~BLEU-4 calculado sobre el subconjunto de patología de MM-Skin.
  El rendimiento sobre imágenes dermoscópicas no fue reportado por separado~\cite{MMSkin2025}.\\
  $^{(\kappa)}$~PeFoMed genera reportes radiológicos (Rx de tórax IU-Xray),
  no dermatológicos; incluido como referencia metodológica para
  la arquitectura VLM~$+$~LoRA en informes médicos~\cite{PeFoMed2024}.
  ROUGE-L y METEOR en escala 0--1 (originalmente 0.286 y 0.157).\\
  La evaluación de \textit{este trabajo} contempla tres niveles:
  métricas automáticas (Nivel~1), coincidencia de entidades clínicas
  mediante NER (Nivel~2) y revisión manual (Nivel~3).
\end{minipage}
\end{table}
Ok, necesitamos hacer una restructuracion de los archivos dentro del projecto. Actualmente tenemos carpetas como `scripts`, `results`, `slurm` o `logs` que contienen muchos archivos sueltos que podrian agruparse por el modelo utilizado. Ten en cuenta que los cambios en estructura de las carpetas tambien se van a tener que ver reflejadas en los distintos scripts que tengamos y en la estructura que utilizamos en los cluster. Cada vez que modifiques la estructura de alguna de estas carpetas, busca donde son utilizadas dentro de los algoritmos y scripts para modificarlas tambien. Con respecto a los logs, actualmente dentro de los clusters tenemos guardados todos los logs de todos los jobs que hemos ejecutado lo que produce que al traer los resultados nos traigan demasiados archivos de logs. Creo que deberiamos modificar los scripts <model>_down.py y agregar una flag separada llamada --logs en el cual se especifique el un numero preciso de logs a descargar o en su defecto un rango de numeros de logs (ej: podemos especificar descargar el log del job N o especificar el rango [N,M]).